\section{Applied Obligation Accounts}

These dossiers test the method without deciding an actual case. Each begins with a person and an act, separates the norms, records the effects each forum can give, and identifies the next responsible step.

\subsection{A professional ordered to participate in disputed conduct}

A licensed employee receives an instruction said to be required by civil regulation and an employment agreement. The employee believes that performing the assigned step would constitute wrongful participation in another act. Four inquiries must remain separate:

\begin{enumerate}
\item The civil lawyer identifies the exact statute, regulation, license condition, employer duty, available accommodation, enforcement authority, and remedy. A supervisor's summary is not yet the governing text.
\item Moral analysis identifies the employee's own chosen act, intention, causal contribution, proximity, alternatives, duress, scandal, and effects on the person served and on coworkers. Labels such as direct, indirect, formal, or material cooperation require the facts that support them.
\item Contract and professional rules determine role duties and procedures for reassignment, notice, referral, continuity, confidentiality, or withdrawal. Those duties may themselves protect vulnerable persons.
\item If performance would require an intrinsically wrongful choice, a well-formed and certain conscience requires refusal. The person should use lawful accommodation and review where possible, prepare for civil or employment consequences, and reduce foreseeable harm to others without performing the prohibited act.
\end{enumerate}

\emph{Evangelium vitae} 73--74 supplies a specific analysis for laws attacking human life, including legislative harm reduction and cooperation (\sourceurl{https://www.vatican.va/content/john-paul-ii/en/encyclicals/documents/hf_jp-ii_enc_25031995_evangelium-vitae.html}{text}). Its concrete conclusions should not be transferred to unrelated disputes without a new moral and legal analysis.

\subsection{A reporting duty and a confidentiality duty}

Suppose a cleric, employee, clinician, teacher, or officer learns facts suggesting abuse or another grave danger. Civil law, canon law, professional regulation, contract, office, and moral duties may all bear on disclosure. The first question is how the information was obtained and what disclosure is proposed.

Information learned in sacramental confession is governed by the inviolable sacramental seal under CIC c.~983 and CCEO c.~733. Information learned in a nonsacramental conversation, administrative file, clinical relationship, or employment investigation requires a different analysis. \emph{Vos estis lux mundi} preserves state-law rights and duties, including reporting duties, while establishing its ecclesiastical reporting framework (\sourceurl{https://www.vatican.va/content/francesco/en/motu_proprio/documents/20230325-motu-proprio-vos-estis-lux-mundi-aggiornato.html}{text}).

The responsible person therefore records the source and capacity in which each fact was received, the exact civil and canonical reporting rule, any privilege or exception, imminent-safety duties, the authorized recipient, required content, deadline, and protections against retaliation. A canonical report, civil report, safeguarding measure, and public statement are different acts. Advice must be obtained without disclosing protected information beyond what the governing rule permits.

\subsection{Marriage in sacramental, canonical, and civil forums}

Two persons can have one history and more than one legally relevant status. Marriage is a natural institution; Christ raised the matrimonial covenant between baptized persons to the dignity of a sacrament (CIC c.~1055). Divine and canonical law govern the marriages of Catholics, while civil authority governs merely civil effects (c.~1059).

Canonical law determines matters including capacity, impediments, consent, form, proof, and tribunal process. Civil law determines civil status, licensing, property, support, custody, and dissolution of civil effects. A civil divorce can alter civil rights without deciding whether a canonical bond arose. A declaration of canonical nullity judges whether the alleged bond validly arose; it does not itself assign every civil consequence.

For the individual, the relevant account includes baptism and Church ascription, prior bonds, domicile, form, dispensations, civil status, dependents, property, pending proceedings, and the relief actually sought. A pastoral judgment, tribunal sentence, and civil judgment possess different authorities and effects. Obligations to support children, avoid scandal, tell the truth, and repair harm can remain while status is contested.

\subsection{Ecclesiastical property held through a civil entity}

A diocesan or religious work may operate through a civil corporation while its assets belong canonically to a public juridic person. A board member can therefore be a civil fiduciary, holder of a canonical office, contractual agent, and moral steward at the same time. Those relations are connected but are not interchangeable.

Natural justice governs honesty, promises, stewardship, fraud, and the common destination of goods. Civil law governs title, corporate authority, contracts, registration, fiduciary duty, creditor protection, and remedies. Canon law governs ecclesiastical ownership and administration and can receive civil contract rules through CIC c.~1290. If an asset is lawfully designated stable patrimony and the applicable thresholds or circumstances are met, the alienation can require valuation, consultation, consent, or permission under cc.~1291--1295.

An alienation of ecclesiastical goods made without the required canonical formalities may be civilly effective yet canonically invalid. When those conditions are present, CIC c.~1296 directs competent ecclesiastical authority to decide whether and how to vindicate Church rights. Before acting, the individual must identify the exact civil and canonical entities, title, asset classification, authority granted, approvals, conflicts of interest, and consequences for beneficiaries and counterparties. Invoking Church ownership does not cure a defective civil instrument; civil signature authority does not supply missing canonical permission.

\subsection{One wrongful act, several forms of responsibility}

An intentional attack on innocent life is contrary to natural moral law and the revealed command against murder. The same event may also fall within civil criminal law, canonical penal law, employment discipline, tort law, and duties of restitution. Each order defines its own elements, subjects, evidence, defenses, and sanctions.

Under current Latin canon law, CIC c.~1321 presumes the accused innocent until the contrary is proved, and punishment requires an externally committed violation gravely imputable by malice or culpability. Once the external violation has been committed, imputability is presumed unless otherwise apparent. Canons 1322--1325 govern excuses and mitigation. Canon 1399 permits punishment for an otherwise unprovided external violation of divine or canon law only when the violation's special gravity demands punishment and urgent need requires prevention or repair of scandal. The CCEO requires analysis under its distinct penal provisions.

Objective moral wrong, moral culpability, civil guilt, canonical imputability, employment cause, and civil liability are therefore separate conclusions. An acquittal for failure of penal proof does not declare an act morally good. A grave moral judgment does not eliminate the prosecution's burden or create a penalty that positive law does not supply.

\subsection{Sunday worship, work, and family duty}

A Latin Catholic is scheduled to work on Sunday and is responsible for dependents. The underlying duties include worship of God, observance of the Lord's Day, support of family, justice to employer and coworkers, and protection of health or public necessity. CIC cc.~1246--1248 specify the Latin obligation to participate in Mass on Sundays and holy days, the period for fulfillment, and the response where participation is impossible because of the absence of a sacred minister or for another grave cause; particular law can alter the holy-day calendar within competence (\sourceurl{https://www.vatican.va/archive/cod-iuris-canonici/eng/documents/cic_lib4-cann1244-1253_en.html}{CIC 1244--1253}).

The individual must establish rite, territory, schedule, available celebrations, the nature and necessity of the work, dependents' needs, contractual options, and any competent dispensation or commutation under c.~1245. A dispensation from an ecclesiastical obligation does not remove the continuing natural and revealed duty to worship God. The natural duty by itself does not select every feast, time, territorial transfer, excuse, or juridical consequence.

\subsection{Religious exercise and public authority}

A person or community seeking to worship, teach, assemble, or organize can invoke a dignity-grounded civil immunity and the applicable constitutional or statutory protection. The civil forum must determine standing, governmental action, the legal standard, evidence, and remedy. \emph{Dignitatis humanae} supplies authoritative Catholic teaching on the natural ground and due limits of the right; it does not itself file a civil claim or determine a state's procedural law.

Public authority may regulate conduct to protect the rights of others and a just public order. The regulation must issue from competent authority, rest on evidence, and be applied fairly; religious error or majority dislike is not enough. The religious actor remains responsible for the rights of neighbors, workers, children, and other affected persons. Natural ground, civil recognition, ecclesiastical mission, and a concrete remedy thus belong to one account without becoming one source.

\subsection{What every dossier establishes}

\begin{studybox}{The practical result}
One act can be morally forbidden, civilly effective, canonically invalid, contractually actionable, nonpenal, and still subject to restitution. Another act can be morally indifferent in itself yet prohibited by a valid and just positive rule, so that compliance---or a competent exception or accommodation---is required. The responsible conclusion states each effect separately, selects the competent forum for relief, and then identifies what this person must do now while protecting the rights of others.
\end{studybox}
